\section{Conclusions}

\subsection{Performance Analysis}

The dual actuator control system successfully demonstrates Varianta C of the laboratory assignment, achieving simultaneous control of a binary actuator (relay) and an analog actuator (PWM) with independent signal conditioning pipelines and unified LCD reporting.

The binary actuator's counter-based debouncing provides reliable switching with a 500\,ms confirmation period (5 cycles at 100\,ms). This is appropriate for relay control, where the mechanical switching time of a typical relay (5--15\,ms) is well within the debounce window. The debounce filter effectively rejects rapid command toggling that could cause relay chatter and premature contact wear.

The analog actuator's conditioning pipeline combines ramping (5\%/cycle rate limit) with the three-stage SignalConditioner (saturation, median, EWMA). The ramping stage limits the rate of duty cycle change to 50\%/second, providing smooth transitions that protect mechanical actuators from sudden load changes. The SignalConditioner further smooths the ramped value, achieving a stable output even with noisy or rapidly changing command inputs.

The FreeRTOS multi-task architecture provides clean separation of concerns: the command task handles serial I/O at the highest priority (50\,ms period), ensuring responsive command processing well within the 100\,ms latency requirement. The actuator tasks run at medium priority (100\,ms period) for timely signal conditioning, while the display task runs at the lowest priority (500\,ms period) since visual updates are less time-critical. The mutex-protected shared state ensures data consistency between tasks without deadlock risk, as all critical sections are short and non-blocking.

\subsection{Limitations}

The current implementation has several limitations that could be addressed in future improvements. The analog actuator is simulated using a PWM-driven LED rather than a real DC motor, so characteristics such as back-EMF, inertia, and current limiting are not represented. The serial command interface is text-based and single-character buffered, which could lose characters under very high input rates. The LCD update rate of 500\,ms is sufficient for human observation but may miss transient states during rapid testing. The system does not implement overload protection or fault detection for the actuators.

\subsection{Improvement Ideas}

Several enhancements could extend the system's capabilities. Adding a motor driver (L298N or L293D) would enable real DC motor control with direction switching and current sensing. Implementing a PID controller for the analog actuator would enable closed-loop speed or position control. Adding EEPROM storage for configuration parameters (debounce count, ramp rate, conditioning parameters) would allow persistent customization. Integrating a keypad as an alternative input device would provide standalone operation without a serial terminal connection.

\subsection{Real-World Impact}

Actuator control with signal conditioning is fundamental to virtually every automated system. The techniques implemented in this lab --- debounced relay switching, ramped PWM control, and multi-task coordination --- are directly applicable to industrial automation (PLC-based production lines), building management systems (HVAC control, lighting automation), automotive systems (window motors, seat positioning, climate control), and consumer electronics (smart home devices, appliance control). The FreeRTOS-based architecture demonstrates the standard approach used in commercial embedded products, where multiple independent control loops must operate concurrently with guaranteed timing and data integrity.
